\chapter{A one-dimensional dust without tangent lines}\label{sec:four-corner-examples}
The four-corner Cantor set is the example following Theorem 3.17 in \cite{Mor16}, fifth edition, pp.~36--37. We prove finite positive length, a quantitative obstruction to tangent lines, and pure unrectifiability. The last step explicitly uses the rectifiable tangent theorem, Theorem \ref{thm:rectifiable-tangents}; the elementary estimates do not claim the source's exact value of the Hausdorff measure.

\begin{definition}[The four-corner set and cylinder measure]\label{def:four-corner-cantor}
Let
\[
 C_4=\left\{\sum_{j\ge1}3\varepsilon_j4^{-j}:\varepsilon_j\in\{0,1\}\right\},\qquad K=C_4\times C_4.
\]
At level $k$, the set $K$ lies in $4^k$ disjoint corner squares of side $4^{-k}$. A cylinder is the intersection of $K$ with one such square. Let $\mu$ be the probability measure assigning mass $4^{-k}$ to each level-$k$ cylinder. Equivalently, push forward Lebesgue probability on $[0,1]^2$ by the two binary-to-base-four digit maps; the dyadic endpoint choices affect a null set only.
\end{definition}

\begin{proposition}[Linear ball bounds and finite positive length]\label{prop:four-corner-linear-measure}
For $x\in K$ and $0<r\le\sqrt2$,
\[
 \frac{r}{8\sqrt2}\le\mu(B_r(x))\le100r.
\]
The upper bound also holds for balls centered outside $K$. Consequently $0<\mathcal H^1(K)<\infty$ and $\dim_{\rm H}K=1$; more precisely $1/100\le\mathcal H^1(K)\le\sqrt2$.
\end{proposition}
\begin{proof}
For $r<1$, choose $k$ with $\ell=4^{-k}\le r<4\ell$. The lower-left corners of level-$k$ squares lie on the grid $\ell\mathbb Z^2$. A square meeting a ball of radius $r$ has each lower-left coordinate in an interval of length at most $2r+\ell<9\ell$. There are at most ten grid choices per coordinate and hence at most one hundred such squares. Each has $\mu$ mass $\ell$, so $\mu(B_r)\le100\ell\le100r$. For $r\ge1$ the bound follows from total mass one.

For the lower bound choose the least $k$ with $\sqrt2\,4^{-k}\le r/2$. The level-$k$ cylinder containing $x$ lies entirely in $B_r(x)$, and minimality gives $4^{-k}>r/(8\sqrt2)$. Thus its mass proves the lower estimate. The same cylinder masses show that every point has zero mass.

The level-$k$ squares cover $K$ with total diameter $4^k\sqrt2\,4^{-k}=\sqrt2$, proving the Hausdorff upper bound. For any sufficiently small countable cover $K\subset\bigcup U_i$, discard sets missing $K$ and choose $x_i\in U_i\cap K$. Then $U_i\cap K\subset\overline B_{\operatorname{diam}U_i}(x_i)$. The ball estimate, obtained for closed balls by enlarging the radius and taking a limit, gives $1\le100\sum_i\operatorname{diam}U_i$. Infimizing and shrinking the cover diameters gives the lower bound. A set of finite positive $\mathcal H^1$ measure has Hausdorff dimension one by the comparison of diameter powers in the Hausdorff definition.
\end{proof}

\begin{proposition}[The cylinder measure is normalized Hausdorff measure]\label{prop:four-corner-hausdorff-measure}
The identity $\mu=\mathcal H^1\llcorner K/\mathcal H^1(K)$ holds. In particular a level-$k$ cylinder has $\mathcal H^1$ measure $4^{-k}\mathcal H^1(K)$.
\end{proposition}
\begin{proof}
Every level-$k$ cylinder is a homothetic copy of $K$ with ratio $4^{-k}$. Hausdorff one-measure scales by that ratio, directly from its definition by diameter covers. Thus normalized Hausdorff measure and $\mu$ agree on every cylinder. Finite unions of cylinders form an algebra, since the partitions refine each other, and this algebra generates the Borel sets of $K$: cylinder diameters tend to zero and every relative neighborhood contains a cylinder around each of its points. Uniqueness of finite measures on a generating algebra gives equality. All measures here are Borel measures on the compact set $K$.
\end{proof}

\begin{lemma}[Differentiation along the Cantor cylinders]\label{lem:four-corner-cylinder-differentiation}
For $f\in L^1(\mu)$, let $Q_k(x)$ be the level-$k$ cylinder containing $x$, and put
\[
 A_kf(x)=\frac1{\mu(Q_k(x))}\int_{Q_k(x)}f\,d\mu.
\]
Then $A_kf(x)\to f(x)$ for $\mu$-almost every $x$. In particular, for a $\mu$-measurable subset $E\subset K$, at almost every point of $E$ the ratio $\mu(Q_k(x)\setminus E)/\mu(Q_k(x))$ tends to zero.
\end{lemma}
\begin{proof}
The cylinder maximal function $Mg=\sup_{k\ge0}A_k|g|$ satisfies
\[
 \mu\{Mg>t\}\le t^{-1}\|g\|_{L^1(\mu)}.
\]
For a finite range of levels, select the first cylinder in each nested chain whose average exceeds $t$. These stopping cylinders are disjoint and cover the set in question; summing $t\mu(Q)<\int_Q|g|$ proves the inequality. Increasing the range gives the full estimate by continuity from below.

For a continuous function on $K$, the averages converge uniformly to the function, because cylinder diameters tend uniformly to zero. Continuous restrictions of compactly supported ambient functions are dense in $L^1(\mu)$ by Theorem \ref{thm:radon-lp-density}. Given such an approximation $h$ to $f$, the limit superior of $|A_kf-f|$ is at most $M(f-h)+|f-h|$. The maximal inequality and Markov's inequality bound the measure where this exceeds $2t$ by $2t^{-1}\|f-h\|_1$. Let the approximation error tend to zero and then take countably many positive rational $t$. This proves almost-everywhere convergence. Apply it to $f=\chi_E$ for the density assertion.
\end{proof}

\begin{lemma}[Every line misses a substantial corner]\label{lem:four-corner-line-obstruction}
Let $Q$ be a level-$k$ square of side $\ell$. For every affine line $L\subset\mathbb R^2$, one of the four child squares $Q'$ satisfies
\[
 \operatorname{dist}(y,L)\ge\frac{\ell}{8\sqrt2}\quad\hbox{for all }y\in Q'.
\]
Its cylinder has $\mathcal H^1$ measure $\ell\mathcal H^1(K)/4$.
\end{lemma}
\begin{proof}
The four child centers form a square of side $3\ell/4$. Project onto a unit normal of $L$. The range of their projections is at least $3\ell/(4\sqrt2)$, so at least one center has distance at least $3\ell/(8\sqrt2)$ from $L$. Every point of its child square is within the half diagonal $\ell/(4\sqrt2)$ of that center. The triangle inequality gives the claimed distance. The measure is the cylinder formula of Proposition \ref{prop:four-corner-hausdorff-measure}.
\end{proof}

\begin{proposition}[No approximate tangent line at any point]\label{prop:four-corner-no-tangent}
For every $x\in K$ and every line $P$ through the origin, $P$ is not an approximate tangent to $K$ at $x$ in Definition \ref{def:approximate-tangent}. More quantitatively, the measure outside a fixed narrow cone about $x+P$ is bounded below by a positive multiple of radius along the cylinder scales.
\end{proposition}
\begin{proof}
Take the parent square of side $\ell=4^{-k}$ containing $x$. Lemma \ref{lem:four-corner-line-obstruction} supplies a child cylinder of mass $\ell\mathcal H^1(K)/4$ at distance at least $\ell/(8\sqrt2)$ from $x+P$. All its points are within $\sqrt2\ell$ of $x$, so their angle from $P$ has sine at least $1/16$. The mass in $B_{2\sqrt2\ell}(x)$ outside the cone of sine aperture $1/32$ is therefore at least $\ell\mathcal H^1(K)/4$.

This also directly contradicts the measure-limit tangent definition. After scaling by $2\sqrt2\ell$, the child lies in the fixed set $\{|y|\le1/2,\operatorname{dist}(y,P)\ge1/32\}$. Choose a continuous compactly supported nonnegative test equal to one on this set and zero on $P$. Its integral against the rescaled measures has a uniform positive lower bound, whereas its integral on $P$ is zero. Hence the required convergence cannot hold.
\end{proof}

\begin{theorem}[Pure unrectifiability of the four-corner set]\label{thm:four-corner-pure-unrectifiability}
For every countably one-rectifiable set $A\subset\mathbb R^2$, $\mathcal H^1(A\cap K)=0$. This is the example in \cite{Mor16}, pp.~36--37. The proof below uses the rectifiable tangent theorem now proved through the local decomposition and cone arguments, Theorem \ref{thm:rectifiable-tangents}, and the locally proved cylinder differentiation lemma.
\end{theorem}
\begin{proof}
It suffices to treat the image of a Lipschitz map on a compact interval, since countably many such images cover any countably rectifiable set up to a null set. Such an image is compact and has finite length, by the elementary diameter-cover bound for a Lipschitz map. Let $E$ be its intersection with $K$ and suppose $\mathcal H^1(E)>0$. Then $E$ is rectifiable, measurable and finite in measure, so Theorem \ref{thm:rectifiable-tangents} gives an approximate tangent at almost every point of $E$. By Lemma \ref{lem:four-corner-cylinder-differentiation} and Proposition \ref{prop:four-corner-hausdorff-measure}, choose such a point $x$ where
\[
 \mathcal H^1((K\cap Q_k(x))\setminus E)=o(4^{-k}).
\]
The child corner furnished by Lemma \ref{lem:four-corner-line-obstruction} for its tangent line contains $\ell\mathcal H^1(K)/4$ of $K$ measure, of which only $o(\ell)$ is missing from $E$. The same continuous test used in Proposition \ref{prop:four-corner-no-tangent} therefore has a uniformly positive integral against the rescaled measures of $E$. This contradicts its approximate tangent. Thus the intersection with every compact Lipschitz curve has zero length; countable union and the null exceptional set prove the assertion.
\end{proof}

\begin{example}[An exceptional projection contains an interval]\label{ex:four-corner-projection-interval}
Both coordinate projections of $K$ have Lebesgue measure zero. In contrast, the orthogonal projection onto the direction $(2,1)/\sqrt5$ is the interval $[0,3/\sqrt5]$.
\end{example}
\begin{proof}
A coordinate projection equals $C_4$, whose level-$k$ cover consists of $2^k$ intervals of length $4^{-k}$; their total length tends to zero. For points of $K$,
\[
 2x+y=3\sum_{j\ge1}(2\varepsilon_j+\delta_j)4^{-j}.
\]
The digit $2\varepsilon_j+\delta_j$ takes each value $0,1,2,3$, independently at every level. These are exactly the base-four expansions of the interval $[0,1]$, including the endpoint expansion with every digit three. Hence $2x+y$ ranges over $[0,3]$. Division by $\sqrt5$ gives the orthogonal projection coordinate. This is a statement about particular directions; it does not purport to prove the almost-everywhere projection theorem for purely unrectifiable sets.
\end{proof}
